% 2021 全国甲卷理科 8：珠峰三角高程测量示意图
% 依据固定上游 PDF 第 1 页；地面投影、竖直高度及仰视线均按原图关系示意。
\begin{tikzpicture}[font=\small,line width=0.45pt,
  x={(0.78cm,0cm)},y={(0cm,0.72cm)},z={(0.32cm,0.20cm)}]
  % 地面投影构成角 A'C'B'=45°, A'B'C'=60° 的三角形（非等比例排版）
  \coordinate (Ap) at (0,0,0);
  \coordinate (Bp) at (4,0,0);
  \coordinate (Cp) at (1.268,0,4.732);
  % 三点位于投影上方；高度仅用于示意
  \coordinate (A) at (0,6.2,0);
  \coordinate (B) at (4,2.464,0);
  \coordinate (C) at (1.268,1,4.732);
  % 地面投影及高度垂线
  \draw (Ap)--(Bp)--(Cp)--cycle;
  \draw[dashed] (Ap)--(A) (Bp)--(B) (Cp)--(C);
  % 测量视线
  \draw (B)--(A);
  \draw (C)--(B);
  % 仰角/地面角文字标注
  \node at (2.0,4.3,0) {$45^\circ$};
  \node at (3.65,2.9,0) {$15^\circ$};
  \node at (2.4,0.8,0) {$45^\circ$};
  \node at (3.6,0.5,0) {$60^\circ$};
  % 高度差标注
  \node[anchor=west] at (5.1,3.2,0) {$BB'-CC'=100$};
  \node[anchor=east] at (-0.1,3.0,0) {$AA'-CC'$};
  % 点名
  \node[anchor=south] at (A) {$A$};
  \node[anchor=west] at (B) {$B$};
  \node[anchor=west] at (C) {$C$};
  \node[anchor=north east] at (Ap) {$A'$};
  \node[anchor=north west] at (Bp) {$B'$};
  \node[anchor=south east] at (Cp) {$C'$};
\end{tikzpicture}
